% §6.2 Sentence-BERT principal component channel, clean 12-city panel,
% with the within-city-centered pooled PCA treatment definition.
% Numbers re-pulled and locked against (refreshed 2026-06-07 after the pooled
% axis was oriented so the pooled estimand is non-negative; PC sign is a
% disclosed convention, commit d4fdeec):
%   results/replications/baur_pooled_pca/baur_pooled_pca_table.csv   (per-city dml_theta/se/ci, oriented)
%   results/replications/baur_pooled_pca/baur_pooled_pca_pooled.json (canonical RE pool, oriented)
%   results/replications/meta_regression_pooled.json (--source baur_pooled): PM+HKSJ theta=+0.156,
%     se=0.042, t(11)=+3.68, two-sided p=0.0036, one-sided 0.0018, CI [+0.063,+0.249], I^2~99%, tau^2_PM~0.021.
% Supersedes the prior v3, which carried the unoriented pool (-0.078) and a
% "Shen-Baur sign disagreement is the central reading" narrative that is an
% artifact of the unidentified principal-component sign (see Section 6.7).

\subsection{Sentence-BERT principal component channel across twelve markets}
\label{sec:baur-12city}

The Baur, Rosenfelder, and Lutz channel \citep{baur2023analysis} treats
listing language as a $768$-dimensional vector from a frozen sentence-BERT
encoder and uses the leading principal component direction as a scalar
treatment in a partially-linear DML estimator. The original study computes the
PCA per market on the local listing distribution. On a single-market analysis
this is the canonical choice, but on a twelve-market panel it produces a
treatment whose direction is identified only locally and only up to a sign, so
per-market PC1 estimates can carry opposite signs across markets even when the
underlying language-price relationship is direction-consistent. A panel built
on per-market axes therefore mixes genuine heterogeneity with an artifact of
local sign and orientation, and it cannot be pooled coherently.

\paragraph{Pooled within-city-centered PCA.} To obtain a treatment definition
that is comparable across markets, we replace the per-market PCAs with a single
global principal-component direction extracted from the pooled within-city
covariance of the stacked embeddings. We center each city's embeddings at its
own mean before pooling, which removes between-city vocabulary differences from
the global covariance and produces a leading eigenvector identified as the
within-city common axis. This is the standard pooling for cross-population PCA
discussed in \citet{flury1984}, equivalent in construction to a linear
discriminant analysis pooled covariance. We do not standardize individual
embedding dimensions before pooling because sentence-BERT outputs are
L2-normalized by design and per-dimension scaling would inflate the anisotropic
noise structure documented by \citet{ethayarajh2019contextual}. We $z$-score the
per-listing projection within each city, so the estimand on each market is a
per-city-standard-deviation move along the common axis. The global axis is
extracted from the within-city-centered pooling of all twelve markets, an
analysis sample of $69{,}173$ listings whose per-market sizes range from $986$ in
San Francisco to $15{,}227$ in New York.

\paragraph{Sign convention.} A principal component and its negation span the
same axis, so the sign of the per-standard-deviation effect is a convention
rather than a finding. We adopt the convention that orients the common axis so
that the pooled effect on log sale price is non-negative, which renders the
per-market estimates sign-comparable to the hedonic uniqueness channel of
Section~\ref{sec:shen-12city} and to the counterfactual differential of
Section~\ref{sec:counterfactual-12}. Under the opposite convention every sign
reported below reverses; the magnitudes, the cross-market dispersion, and the
concordance with the other channels are invariant to it. What the orientation
does not do is manufacture agreement between channels, a point we develop in
Section~\ref{sec:cross-method-12}, because the across-market correlation between
the BERT axis and the hedonic channel is invariant in magnitude under any global
sign flip of either one.

\paragraph{Per-market estimates.} Table~\ref{tab:baur-12city} reports the
per-market per-standard-deviation effect under the oriented pooled-PCA
treatment. The corrected definition removes the sign ambiguity of the per-market
axes and delivers a single direction of effect across the panel. Every one of
the twelve markets returns a confidence interval that excludes zero at the
$95\%$ level, and eleven of the twelve carry a positive point estimate, with New
York the sole exception at $-0.019$ on an interval $[-0.033, -0.004]$ that is
small and barely separated from zero. Philadelphia and Chicago anchor the high
end at $+0.411$ on $[+0.390, +0.432]$ and $+0.456$ on $[+0.419, +0.493]$,
followed by Atlanta at $+0.243$ and Dallas at $+0.195$. San Francisco returns
$+0.101$ on $[+0.051, +0.151]$, the District of Columbia $+0.118$, and Denver
$+0.103$, while Boston, Portland, Phoenix, and Seattle return the smallest
positive effects, $+0.072$, $+0.082$, $+0.056$, and $+0.049$. Portland, which
the per-market PCA had rendered a degenerate outlier, now sits squarely inside
the panel at $+0.082$ on $[+0.057, +0.107]$, so the panel is clean at twelve
markets and there is no market we exclude. The per-market intervals are tight
because the pooled-PCA treatment is estimated on the full within-market sample,
which ranges into the thousands of listings, rather than on the description
subset.

\begin{table}[t]
\centering
\caption{Sentence-BERT pooled-PCA channel: per-market DML effect per standard
deviation along the oriented common axis on log sale price, with
influence-function $95\%$ intervals, ordered by effect size. The axis is oriented
so the pooled effect is non-negative; the principal-component sign is a disclosed
convention (text). All twelve intervals exclude zero.}
\label{tab:baur-12city}
\small
\begin{tabular}{lrrc}
\toprule
Market & $n$ & $\hat\theta$ & $95\%$ CI \\
\midrule
Chicago       & $5{,}576$  & $+0.456$ & $[+0.419, +0.493]$ \\
Philadelphia  & $7{,}030$  & $+0.411$ & $[+0.390, +0.432]$ \\
Atlanta       & $6{,}134$  & $+0.243$ & $[+0.220, +0.266]$ \\
Dallas        & $8{,}008$  & $+0.195$ & $[+0.175, +0.215]$ \\
Washington    & $4{,}015$  & $+0.118$ & $[+0.087, +0.149]$ \\
Denver        & $5{,}280$  & $+0.103$ & $[+0.084, +0.121]$ \\
San Francisco & $986$      & $+0.101$ & $[+0.051, +0.151]$ \\
Portland      & $4{,}337$  & $+0.082$ & $[+0.057, +0.107]$ \\
Boston        & $2{,}610$  & $+0.072$ & $[+0.029, +0.115]$ \\
Phoenix       & $7{,}086$  & $+0.056$ & $[+0.030, +0.082]$ \\
Seattle       & $2{,}884$  & $+0.049$ & $[+0.019, +0.080]$ \\
New York      & $15{,}227$ & $-0.019$ & $[-0.033, -0.004]$ \\
\midrule
Pooled (PM, HKSJ) & & $+0.156$ & $[+0.063, +0.249]$ \\
\bottomrule
\end{tabular}
\end{table}

\paragraph{Pooled inference.} Section~\ref{sec:meta-regression-12} takes up the
pooled inference question directly via random-effects meta-analysis of the
twelve per-market estimates. Combining them with a Paule-Mandel $\tau^2$ and the
Hartung-Knapp-Sidik-Jonkman variance correction returns a pooled
per-standard-deviation effect of $+0.156$ with standard error $0.042$, a $t$
statistic of $+3.68$ on $11$ degrees of freedom, a two-sided $p$ of $0.004$, and
a $95\%$ confidence interval of $[+0.063, +0.249]$. The pooled estimate is
detectable and positive across the panel under the adopted orientation.
Heterogeneity is large, with an $I^2$ near $99\%$ and a Paule-Mandel $\tau^2$ of
about $0.021$, so the pooled point estimate masks substantial cross-market
dispersion that Chicago and Philadelphia anchor at one end and New York at the
other. Because between-market variance dwarfs the within-market sampling
variance, the meta-analytic weights are close to uniform and the pooled estimate
sits near the unweighted across-market mean. We attempted a single pooled DML
estimator with city fixed effects and the Chiang-Kato-Ma-Sasaki
\citep{chiang2022multiway} cluster-robust variance, but the small number of
clusters produced fold-assignment instability in the cross-fitted nuisances and
a Dirichlet cluster bootstrap that did not stabilize, with pooled point estimates
that shifted materially across independent invocations on the same seed and data,
so we report it only as a robustness check and not as the headline. The
random-effects meta-analysis treats each market's $\hat\theta_i$ as a sufficient
statistic and combines them via inverse-variance weighting with a Paule-Mandel
$\tau^2$, which is the canonical small-$G$ alternative to cluster-respecting
cross-fitting \citep{veroniki2016methods}.

\paragraph{Magnitude and interpretation.} The pooled per-standard-deviation
effect of $+0.156$ on log sale price is about a fifteen percent change in price,
which is small as a fraction of value but substantial against the magnitudes
published in the housing and text-as-treatment literatures. A single additional
bedroom in the canonical \citet{sirmans2005composition} meta-analysis of
forty-five hedonic studies carries a mean effect of $+3.8\%$ on log price, so the
pooled text-axis effect is roughly four times the marginal price of an additional
bedroom in magnitude, and it is several times the racial sale-price gap of about
three percent that \citet{bayer2017racial} document for Black and Hispanic
homebuyers across Chicago, Baltimore-Washington, San Francisco, and Los Angeles.
The more informative comparison is internal. The Doc2Vec uniqueness construct of
\citet{shen2021information} runs in the same direction as the oriented BERT axis,
both within markets and across the panel. In San Francisco, where both channels
are estimated on the identical sample, the same DML pipeline recovers a
uniqueness effect of $+0.122$ and a BERT-PC1 effect of $+0.101$, and pooled
across the twelve markets the uniqueness effect is about $+0.07$ against the
$+0.156$ of the BERT axis. The two text treatments therefore concord in sign and
in cross-market ranking, with the BERT axis roughly twice the hedonic channel in
pooled magnitude. We are deliberate about what this concordance does and does not
establish. The two channels are alternative encodings of the same listing text
applied to the same listings under the same confounder adjustment, and
Section~\ref{sec:cross-method-12} shows they are nearly collinear across markets,
so their agreement is evidence that the text-on-price signal is robust to the
choice of text representation rather than evidence from two independent sources.
Textual distinctiveness in the Shen sense and the utilitarian-to-amenity axis
that the leading sentence-BERT component recovers price in the same direction,
and the difference between them is one of magnitude and of which markets each
detects, not of sign.
